% Abkürzungen, ggfs. mit korrektem Leerraum 
\newcommand{\bs}{$\backslash$\xspace}
\newcommand{\bspw}{bspw.\xspace}
\newcommand{\bzw}{bzw.\xspace}
\newcommand{\ca}{ca.\xspace}
\newcommand{\dahe}{\mbox{d.\,h.}\xspace}
\newcommand{\etc}{etc.\xspace}
\newcommand{\eur}[1]{\mbox{#1\,\texteuro}\xspace}
\newcommand{\evtl}{evtl.\xspace}
\newcommand{\ggfs}{ggf.\xspace}
\newcommand{\Ggfs}{Ggf.\xspace}
\newcommand{\gqq}[1]{\glqq{}#1\grqq{}}
\newcommand{\inkl}{inkl.\xspace}
\newcommand{\insb}{insb.\xspace}
\newcommand{\ua}{\mbox{u.\,a.}\xspace}
\newcommand{\usw}{usw.\xspace}
\newcommand{\Vgl}{Vgl.\xspace}
\newcommand{\zB}{\mbox{z.\,B.}\xspace}
\newcommand{\og}{\mbox{o.\,g.}\xspace}
\newcommand{\sog}{sog.\xspace}

% Befehle für häufig anfallende Aufgaben
\newcommand{\Abbildung}[1]{Abbildung~\ref{#1}}
\newcommand{\AbbildungV}[1]{Abbildung~\ref{#1}: \nameref{#1} auf S.~\pageref{#1}}
\newcommand{\Anhang}[1]{Anhang~\ref{#1}: \nameref{#1} auf S.~\pageref{#1}}
\newcommand{\AnhangK}[1]{Anhang~\ref{#1} auf S.~\pageref{#1}}
\newcommand{\AnhangKapitel}[1]{Kapitel~\ref{#1}: \nameref{#1} auf S.~\pageref{#1}}
\newcommand{\Tabelle}[1]{Tabelle~\ref{#1}: \nameref{#1}}
\newcommand{\TabelleK}[1]{Tabelle~\ref{#1} auf S.~\pageref{#1}}
\newcommand{\TabelleV}[1]{Tabelle~\ref{#1}: \nameref{#1} auf S.~\pageref{#1}}

\newcommand{\includegraphicsKeepAspectRatio}[2]{\includegraphics[width=#2\textwidth,height=#2\textheight,keepaspectratio]{#1}}
\newcommand{\Zitat}[2][\empty]{\ifthenelse{\equal{#1}{\empty}}{\citep{#2}}{\citep[#1]{#2}}}
\newcommand{\Autor}[1]{\textsc{#1}} % zum Ausgeben von Autoren
\newcommand{\itemd}[2]{\item{\textbf{#1}}\\{#2}} % erzeugt ein Listenelement mit fetter Überschrift

% fügt Tabellen aus einer TEX-Datei ein
\newcommand{\newTabelle}[3] % Parameter: caption, label, file
{\begin{table}[htb]
\centering
\singlespacing
\input{Tabellen/#3}
\caption{#1}
\label{#2}
\end{table}}

\newcommand{\newTabelleAnhang}[1] % Parameter: file
{\begin{center}
\singlespacing
\input{Tabellen/#1}
\end{center}}

% einfaches Wechseln der Schrift, z.B.: \changefont{cmss}{sbc}{n}
\newcommand{\changefont}[3]{\fontfamily{#1} \fontseries{#2} \fontshape{#3} \selectfont}

% Verwendung analog zu \includegraphics
\newlength{\myx} % Variable zum Speichern der Bildbreite
\newlength{\myy} % Variable zum Speichern der Bildhöhe
\newcommand\includegraphicstotab[2][\relax]{%
% Abspeichern der Bildabmessungen
\settowidth{\myx}{\includegraphics[{#1}]{#2}}%
\settoheight{\myy}{\includegraphics[{#1}]{#2}}%
% das eigentliche Einfügen
\parbox[c][1.1\myy][c]{\myx}{%
\includegraphics[{#1}]{#2}}%
}



% Tabellen, die den Zwischenstand nach einer Projektphase verdeutlichen
\newcommand{\Zwischenstand}[2]{\subsection{Zwischenstand}Tabelle~\ref{tab:#2} zeigt den Zwischenstand nach der #1.\newTabelle{Zwischenstand nach der #1}{tab:#2}{#2.tex}}

% Verweise
\newcommand{\Verweis}[1]{Kapitel~\ref{#1}}
\newcommand{\VerweisKL}[1]{Kapitel~\ref{#1} auf Seite \pageref{#1}}
\newcommand{\VerweisVL}[1]{Kapitel~\ref{#1}: \nameref{#1} auf Seite \pageref{#1}}
\newcommand{\VerweisKT}[1]{(Siehe Kapitel~\ref{#1} auf Seite \pageref{#1})}
\newcommand{\VerweisVT}[1]{(Siehe Kapitel~\ref{#1}: \nameref{#1} auf Seite \pageref{#1})}
\newcommand{\VerweisKF}[1]{\footnote{Siehe Kapitel~\ref{#1} auf Seite \pageref{#1}.}}
\newcommand{\VerweisVF}[1]{\footnote{Siehe Kapitel~\ref{#1}: \nameref{#1} auf Seite \pageref{#1}.}}

% ### Semantische Textverarbeitung ###

% Abkürzung
\newcommand{\Abkuerzung}[1]{\textit{\ac{#1}}}

% Begriffe
\newcommand{\Fachbegriff}[1]{\textbf{#1}}
\newcommand{\Klasse}[1]{\texttt{#1}}
\newcommand{\Methode}[1]{\texttt{#1}}
\newcommand{\MethodenRef}[1]{\texttt{#1}}
\newcommand{\Code}[1]{\textsf{#1}}
